\subsection{Morfismos planos de espacios anillados}
\label{subsection:0.6.7-fr}

\begin{env}[6.7.1]
\label{0.6.7.1-fr}
Sea $f:X\to Y$ un morfismo de espacios anillados, y sea $\mathcal F$ un
$\mathcal O_X$-Módulo. Se dice que $\mathcal F$ es \emph{$f$-plano} (o $Y$-\emph{plano}
cuando no hay que temer ninguna confusión acerca de $f$) \emph{en un punto} $x\in X$ si
$\mathcal F_x$ es un $\mathcal O_{f(x)}$-módulo plano; se dice que
$\mathcal F$ es \emph{$f$-plano sobre} $y\in Y$ si $\mathcal F$ es $f$-plano
en todos los puntos $x\in f^{-1}(y)$; se dice que $\mathcal F$ es \emph{$f$-plano}
si $\mathcal F$ es $f$-plano en todos los puntos de $X$. Se dice que el
morfismo $f$ es \emph{plano en} $x\in X$ (resp. \emph{plano sobre} $y\in Y$, resp.
\emph{plano}) si $\mathcal O_X$ es $f$-plano en $x$ (resp. $f$-plano sobre $y$,
resp. $f$-plano).
\end{env}

\begin{env}[6.7.2]
\label{0.6.7.2-fr}
Con las notaciones de (6.7.1), si $\mathcal F$ es $f$-plano en $x$, para toda
vecindad abierta $U$ de $y=f(x)$, el funtor
$(f^*(\mathcal G)\otimes_{\mathcal O_X}\mathcal F)_x$ en $\mathcal G$ es
\emph{exacto} en la categoría de los $(\mathcal O_Y|U)$-Módulos; en efecto, este
tallo se identifica canónicamente con
$\mathcal G_y\otimes_{\mathcal O_y}\mathcal F_x$, y nuestra afirmación resulta
de la definición. En particular, si $f$ es un morfismo \emph{plano}, el funtor
$f^*$ es \emph{exacto} en la categoría de los $\mathcal O_Y$-Módulos.
\end{env}

\begin{env}[6.7.3]
\label{0.6.7.3-fr}
Recíprocamente, supongamos que el haz de anillos $\mathcal O_Y$ es \emph{coherente}, y
supongamos que para \emph{toda} vecindad abierta $U$ de $y$, el funtor
$(f^*(\mathcal G)\otimes_{\mathcal O_X}\mathcal F)_x$ es exacto en
$\mathcal G$ en la categoría de los $(\mathcal O_Y|U)$-Módulos \emph{coherentes}.
Entonces $\mathcal F$ es \emph{$f$-plano en $x$}. En efecto, basta probar que
para todo ideal de tipo finito $\mathfrak J$ de $\mathcal O_y$,
el homomorfismo canónico
$\mathfrak J\otimes_{\mathcal O_y}\mathcal F_x\to\mathcal F_x$ es inyectivo
(6.1.1). Ahora bien, se sabe (5.3.8) que existe entonces una vecindad

\oldpage[0\textsubscript{I}]{60}

abierta $U$ de $y$ y un haz coherente de ideales $\mathcal J$ de
$\mathcal O_Y|U$ tal que $\mathcal J_y=\mathfrak J$, de donde la conclusión.
\end{env}

\begin{env}[6.7.4]
\label{0.6.7.4-fr}
Los resultados de (6.1) sobre los módulos planos se transcriben inmediatamente en
proposiciones sobre los haces $f$-planos en un punto:

Si $0\to\mathcal F'\to\mathcal F\to\mathcal F''\to0$ es una sucesión exacta
de $\mathcal O_X$-Módulos y si $\mathcal F''$ es $f$-plano en el punto
$x\in X$, entonces, para toda vecindad abierta $U$ de $y=f(x)$ y todo
$(\mathcal O_Y|U)$-Módulo $\mathcal G$, la sucesión
\[
  0\longrightarrow
  (f^*(\mathcal G)\otimes_{\mathcal O_X}\mathcal F')_x
  \longrightarrow
  (f^*(\mathcal G)\otimes_{\mathcal O_X}\mathcal F)_x
  \longrightarrow
  (f^*(\mathcal G)\otimes_{\mathcal O_X}\mathcal F'')_x
  \longrightarrow0
\]
es exacta. Para que $\mathcal F$ sea $f$-plano en $x$, es entonces necesario y suficiente
que $\mathcal F'$ lo sea. Se tienen enunciados análogos para las nociones
correspondientes de $\mathcal O_X$-Módulo $f$-plano sobre $y\in Y$, o
de $\mathcal O_X$-Módulo $f$-plano.
\end{env}

\begin{env}[6.7.5]
\label{0.6.7.5-fr}
Sean $f:X\to Y$, $g:Y\to Z$ dos morfismos de espacios anillados; sean
$x\in X$, $y=f(x)$, $\mathcal F$ un $\mathcal O_X$-Módulo. Si $\mathcal F$
es $f$-plano en el punto $x$ y si el morfismo $g$ es plano en el punto $y$, entonces
$\mathcal F$ es $(g\circ f)$-plano en $x$ (6.2.1). En particular, si $f$ y
$g$ son morfismos planos, $g\circ f$ es plano.
\end{env}

\begin{env}[6.7.6]
\label{0.6.7.6-fr}
Sean $X$, $Y$ dos espacios anillados, $f:X\to Y$ un morfismo \emph{plano}. Entonces
el homomorfismo canónico de bifuntores (4.4.6)
\[
  f^*(\mathcal{H}om_{\mathcal O_Y}(\mathcal F,\mathcal G))
  \longrightarrow
  \mathcal{H}om_{\mathcal O_X}(f^*(\mathcal F),f^*(\mathcal G))
  \tag{6.7.6.1}
  \label{0.6.7.6.1-fr}
\]
es un \emph{isomorfismo} cuando $\mathcal F$ admite una \emph{presentación finita}
(5.2.5).

En efecto, siendo la cuestión local, se puede suponer que existe una sucesión
exacta $\mathcal O_Y^m\to\mathcal O_Y^n\to\mathcal F\to0$. Ahora bien, los dos
miembros de (6.7.6.1) son funtores exactos por la izquierda en $\mathcal F$ en
virtud de la hipótesis sobre $f$; se reduce entonces a demostrar la proposición
cuando $\mathcal F=\mathcal O_Y$, caso en que es trivial.
\end{env}

\begin{env}[6.7.8]
\label{0.6.7.8-fr}
Se dice que un morfismo $f:X\to Y$ de espacios anillados es \emph{fielmente plano} si
$f$ es \emph{sobreyectivo} y si, para todo $x\in X$, $\mathcal O_x$ es un
$\mathcal O_{f(x)}$-módulo \emph{fielmente plano}. Cuando $X$ y $Y$ son
espacios localmente anillados (5.5.1), equivale a decir que el
morfismo $f$ es \emph{sobreyectivo} y \emph{plano} (6.6.2). Cuando $f$ es fielmente plano,
$f^*$ es un funtor \emph{exacto} y \emph{fiel} en la categoría de los
$\mathcal O_Y$-Módulos (6.6.1, $a)$) y para que un $\mathcal O_Y$-Módulo
$\mathcal G$ sea $Y$-plano, es necesario y suficiente que $f^*(\mathcal G)$ lo
sea (6.6.3).
\end{env}

